% Isosyll v4.0.0 — Executive Brief (CEOs, product leaders & operations executives)
% Strict 2-page US Letter PDF. Every measured number is a macro from the paper's generated
% docs/paper/data/metrics.tex; heatmap cell colours come from docs/paper/data/cold_dev.csv.
\documentclass[letterpaper]{article}
\pdfminorversion=7
\usepackage{isosyll-brief-style}
% Generated by python/ultra/tools/paper_data.py. Do not edit by hand.
% Every macro is a bare number (no unit, no %). Provenance: docs/paper/data/SOURCES.md.
% Labels: exploratory = shared-host / non-qualifying measurement; reported = configuration
% or descriptive constant; qualified = from a qualifying window (none yet).
\providecommand{\PaperMissing}[1]{\textcolor{red}{??}}

\newcommand{\AASpread}{25}% [exploratory] % CPU
\newcommand{\AbabBlocks}{8}% [reported] blocks
\newcommand{\AbabPholmDense}{\ensuremath{1.4\times10^{-10}}}% [exploratory] p
\newcommand{\AbabPholmMR}{\ensuremath{1.9\times10^{-6}}}% [exploratory] p
\newcommand{\AbabRatioDense}{0.666}% [exploratory] x
\newcommand{\AbabRatioDenseHi}{0.678}% [exploratory] x
\newcommand{\AbabRatioDenseLo}{0.655}% [exploratory] x
\newcommand{\AbabRatioMR}{0.629}% [exploratory] x
\newcommand{\AbabRatioMRHi}{0.684}% [exploratory] x
\newcommand{\AbabRatioMRLo}{0.578}% [exploratory] x
\newcommand{\AbabSpeedDense}{1.50}% [exploratory] x
\newcommand{\AbabSpeedMR}{1.59}% [exploratory] x
\newcommand{\AvgGainClean}{0.43}% [exploratory] pp WER
\newcommand{\AvgGainOther}{0.84}% [exploratory] pp WER
\newcommand{\BOneAttnCut}{3--10}% [exploratory] %
\newcommand{\BeamGainClean}{0.63}% [exploratory] pp WER
\newcommand{\BeamGainOther}{1.19}% [exploratory] pp WER
\newcommand{\BridgeStep}{5{,}000}% [reported] steps
\newcommand{\CCXCount}{16}% [reported] CCX
\newcommand{\ChunkSixteenGainClean}{0.39}% [exploratory] pp WER
\newcommand{\ChunkSixteenGainOther}{0.90}% [exploratory] pp WER
\newcommand{\ColdRowsCount}{0}% [reported] rows
\newcommand{\CoresPerCCX}{4}% [reported] cores
\newcommand{\CostClaimDense}{0.0369}% [exploratory] CPU-s/audio-s
\newcommand{\CostClaimMR}{0.0279}% [exploratory] CPU-s/audio-s
\newcommand{\CostClaimRzero}{0.0864}% [exploratory] CPU-s/audio-s
\newcommand{\CostEngineFP}{0.104}% [exploratory] CPU-s/audio-s
\newcommand{\CostEngineIntEight}{0.041}% [exploratory] CPU-s/audio-s
\newcommand{\CostLadderDense}{0.0443}% [exploratory] CPU-s/audio-s
\newcommand{\CostLadderMR}{0.0354}% [exploratory] CPU-s/audio-s
\newcommand{\CostLadderRzero}{0.0864}% [exploratory] CPU-s/audio-s
\newcommand{\CostNemo}{0.472}% [exploratory] CPU-s/audio-s
\newcommand{\CostORT}{0.175}% [exploratory] CPU-s/audio-s
\newcommand{\CostRatioDense}{1.95}% [exploratory] x
\newcommand{\CostRatioMR}{2.44}% [exploratory] x
\newcommand{\CostRzero}{0.078}% [exploratory] CPU-s/audio-s
\newcommand{\DenDenseCCX}{96}% [exploratory] streams/CCX
\newcommand{\DenDenseCore}{24}% [exploratory] streams/core
\newcommand{\DenFailDenseCCX}{112}% [exploratory] streams/CCX
\newcommand{\DenFailMRCCX}{160}% [exploratory] streams/CCX
\newcommand{\DenFailRzeroCCX}{48}% [exploratory] streams/CCX
\newcommand{\DenMRCCX}{144}% [exploratory] streams/CCX
\newcommand{\DenMRCore}{36}% [exploratory] streams/core
\newcommand{\DenPninetyNineDense}{596}% [exploratory] ms
\newcommand{\DenPninetyNineMR}{735}% [exploratory] ms
\newcommand{\DenPninetyNineRzero}{729}% [exploratory] ms
\newcommand{\DenRatioDense}{3.0}% [exploratory] x
\newcommand{\DenRatioMR}{4.5}% [exploratory] x
\newcommand{\DenRatioMRvsDense}{1.5}% [exploratory] x
\newcommand{\DenRungSeconds}{20}% [reported] s
\newcommand{\DenRzeroCCX}{32}% [exploratory] streams/CCX
\newcommand{\DenRzeroCore}{8}% [exploratory] streams/core
\newcommand{\DenWorkers}{4}% [reported] workers
\newcommand{\DtwoAudioHoursRehearsal}{6303}% [reported] audio-h
\newcommand{\DtwoEpoch}{8}% [reported] epoch (0-indexed)
\newcommand{\DtwoEpochRehearsal}{6}% [reported] epoch (0-indexed)
\newcommand{\DtwoStep}{60{,}000}% [reported] steps
\newcommand{\EngineBeamMRDelta}{\ensuremath{-}1.43}% [exploratory] pp WER
\newcommand{\FrontierDense}{0.0435}% [exploratory] CPU-s/audio-s
\newcommand{\FrontierMR}{0.0330}% [exploratory] CPU-s/audio-s
\newcommand{\FrontierPerCoreDense}{23.0}% [exploratory] streams/core
\newcommand{\FrontierPerCoreMR}{30.3}% [exploratory] streams/core
\newcommand{\FrontierPerCoreRzero}{14.3}% [exploratory] streams/core
\newcommand{\FrontierRzero}{0.0698}% [exploratory] CPU-s/audio-s
\newcommand{\FrontierStreams}{16}% [reported] streams
\newcommand{\FullAttnLossClean}{0.81}% [exploratory] pp WER
\newcommand{\GateBeamDelta}{+0.007}% [exploratory] pp WER
\newcommand{\GateBeamUpper}{+0.150}% [exploratory] pp WER
\newcommand{\GateDenseDelta}{+0.021}% [exploratory] pp WER
\newcommand{\GateDenseUpper}{+0.157}% [exploratory] pp WER
\newcommand{\GateDtwoAvgDelta}{+0.042}% [exploratory] pp WER
\newcommand{\GateDtwoAvgUpper}{+0.140}% [exploratory] pp WER
\newcommand{\GateMRDelta}{+0.069}% [exploratory] pp WER
\newcommand{\GateMRUpper}{+0.203}% [exploratory] pp WER
\newcommand{\GateMargin}{0.10}% [reported] pp WER
\newcommand{\GateMaxUpper}{+0.211}% [exploratory] pp WER
\newcommand{\GemmShare}{75--80}% [exploratory] %
\newcommand{\GreedyDecoderCost}{0.0016}% [exploratory] CPU-s/audio-s
\newcommand{\IntEightMRDenCCX}{136}% [exploratory] streams/CCX
\newcommand{\IntEightUncentred}{+0.207}% [exploratory] pp WER
\newcommand{\LThreeMiB}{16}% [reported] MiB
\newcommand{\LeftSixtyFourAttn}{0.70}% [exploratory] x of 128
\newcommand{\LeftSixtyFourDelta}{+0.09}% [exploratory] pp WER
\newcommand{\LeftThirtyTwoAttn}{0.45}% [exploratory] x of 128
\newcommand{\LeftThirtyTwoDelta}{+0.68}% [exploratory] pp WER
\newcommand{\MRBeamPanelBeamFP}{14.43}% [exploratory] % WER
\newcommand{\MRBeamPanelBeamInt}{14.44}% [exploratory] % WER
\newcommand{\MRBeamPanelGreedyFP}{15.86}% [exploratory] % WER
\newcommand{\MRPadDeltaMax}{0.03}% [exploratory] pp WER
\newcommand{\MatchedRung}{32}% [reported] streams/CCX
\newcommand{\NemoLongCost}{0.451}% [exploratory] CPU-s/audio-s
\newcommand{\NemoLongTestClean}{2.58}% [exploratory] % WER
\newcommand{\NemoLongTestOther}{5.18}% [exploratory] % WER
\newcommand{\NemoShortCost}{0.709}% [exploratory] CPU-s/audio-s
\newcommand{\NemoShortTestClean}{2.65}% [exploratory] % WER
\newcommand{\NemoShortTestOther}{5.60}% [exploratory] % WER
\newcommand{\PanelSpeakers}{73}% [reported] speakers
\newcommand{\ParamsDtwo}{58.7}% [reported] M parameters
\newcommand{\ParamsDtwoMR}{59.2}% [reported] M parameters
\newcommand{\ParityFP}{\ensuremath{2.4\times10^{-6}}}% [exploratory] abs error
\newcommand{\QkvCapRatio}{1.5}% [reported] x
\newcommand{\QkvRunawayFrom}{70}% [exploratory] norm
\newcommand{\QkvRunawayTo}{684}% [exploratory] norm
\newcommand{\RssEngineFP}{478}% [exploratory] MB
\newcommand{\RssEngineIntEight}{646}% [exploratory] MB
\newcommand{\RssEngineIntEightPrepared}{259}% [reported] MB
\newcommand{\RssNemo}{1{,}284}% [exploratory] MB
\newcommand{\RssORT}{776}% [exploratory] MB
\newcommand{\RssRzero}{239}% [exploratory] MB
\newcommand{\SingleMthreeDense}{0.0132}% [exploratory] CPU-s/audio-s
\newcommand{\SingleMthreeMR}{0.0123}% [exploratory] CPU-s/audio-s
\newcommand{\SingleZenDense}{0.040}% [exploratory] CPU-s/audio-s
\newcommand{\SpeedupEngineFPVsORT}{1.7}% [exploratory] x
\newcommand{\SpeedupVsORT}{4.2}% [exploratory] x
\newcommand{\SpeedupVsRzero}{1.9}% [exploratory] x
\newcommand{\StackGainClean}{0.99}% [exploratory] pp WER
\newcommand{\StackGainOther}{2.02}% [exploratory] pp WER
\newcommand{\StreamsPerCoreMthreeMR}{84}% [exploratory] streams/core
\newcommand{\StreamsPerCoreMthreeMRFail}{96}% [exploratory] streams/core
\newcommand{\SweepDtwoBeamFourPooled}{10.82}% [exploratory] % WER
\newcommand{\SweepDtwoBestBeam}{8}% [exploratory] beam
\newcommand{\SweepDtwoBestPooled}{10.69}% [exploratory] % WER
\newcommand{\SweepDtwoGreedyPooled}{11.67}% [exploratory] % WER
\newcommand{\SweepDtwoMRBeamFourPooled}{14.41}% [exploratory] % WER
\newcommand{\SweepDtwoMRBestBeam}{8}% [exploratory] beam
\newcommand{\SweepDtwoMRBestPooled}{14.20}% [exploratory] % WER
\newcommand{\SweepDtwoMRGreedyPooled}{15.89}% [exploratory] % WER
\newcommand{\TernaryDenCCX}{112}% [exploratory] streams/CCX
\newcommand{\TernaryEngineWer}{17.98}% [exploratory] % WER
\newcommand{\TernaryMiB}{14.26}% [reported] MiB
\newcommand{\TernaryZeros}{31.8}% [reported] %
\newcommand{\TrainHours}{960}% [reported] h
\newcommand{\UsevenBar}{1.2}% [reported] x
\newcommand{\UsevenVsActHi}{0.92}% [exploratory] x
\newcommand{\UsevenVsActLo}{0.89}% [exploratory] x
\newcommand{\UsevenVsKzeroHi}{1.01}% [exploratory] x
\newcommand{\UsevenVsKzeroLo}{0.94}% [exploratory] x
\newcommand{\WerBridgeFP}{21.88}% [exploratory] % WER
\newcommand{\WerBridgeIntEight}{21.80}% [exploratory] % WER
\newcommand{\WerDtwoAvgGreedyClean}{7.22}% [exploratory] % WER
\newcommand{\WerDtwoAvgGreedyOther}{17.73}% [exploratory] % WER
\newcommand{\WerDtwoBeamClean}{7.03}% [exploratory] % WER
\newcommand{\WerDtwoBeamOther}{17.39}% [exploratory] % WER
\newcommand{\WerDtwoBestClean}{6.66}% [exploratory] % WER
\newcommand{\WerDtwoBestOther}{16.55}% [exploratory] % WER
\newcommand{\WerDtwoBestPanelClean}{6.93}% [exploratory] % WER
\newcommand{\WerDtwoBestPanelOther}{16.08}% [exploratory] % WER
\newcommand{\WerDtwoGreedyClean}{7.65}% [exploratory] % WER
\newcommand{\WerDtwoGreedyOther}{18.57}% [exploratory] % WER
\newcommand{\WerDtwoMRClean}{11.10}% [exploratory] % WER
\newcommand{\WerDtwoMROther}{22.90}% [exploratory] % WER
\newcommand{\WerLadderDenseDevClean}{8.95}% [exploratory] % WER
\newcommand{\WerLadderDensePanel}{14.19}% [exploratory] % WER
\newcommand{\WerLadderMRDevClean}{11.19}% [exploratory] % WER
\newcommand{\WerLadderMRPanel}{15.93}% [exploratory] % WER
\newcommand{\WerLadderMRPanelClean}{11.56}% [exploratory] % WER
\newcommand{\WerNemoDevClean}{2.16}% [exploratory] % WER
\newcommand{\WerRzeroDevClean}{2.85}% [exploratory] % WER
\newcommand{\XisaAgree}{7{,}300}% [exploratory] utterances
\newcommand{\XisaModels}{5}% [reported] checkpoints
\newcommand{\XisaPanel}{730}% [reported] utterances
\newcommand{\XisaRfourPre}{636}% [exploratory] of 730
\newcommand{\XisaTotal}{7{,}300}% [exploratory] utterances

\newcommand{\RehearsalCaption}{Warm single-harness rehearsal (exploratory): R0 Zipformer INT8, D2 step 45k dense INT8 and D2MR step 50k multirate INT8 in one \texttt{ultra bench-suite} binary on a shared AMD Threadripper PRO 3995WX (Zen~2), \DenWorkers{} single-threaded workers on one CCX, \DenRungSeconds{}~s ladder rungs. Left: whole-service CPU-seconds per audio-second at the matched \MatchedRung-stream rung against the DEV400 dev-clean WER of the same checkpoints. Right: highest sustained streams per CCX; the WER of each checkpoint is printed on its bar.}
\newcommand{\ColdCaption}{Cold process-tree CPU cost against pooled DEV WER on the frozen 730-utterance DEV panel, one pinned physical core, median of three repetitions, AMD Threadripper PRO 3995WX (Zen~2).}

% Generated by docs/briefs/scripts/gen_brief_heatmap.py from docs/paper/data/cold_dev.csv.
% Do not edit by hand. Pastel-to-rich perceptual heatmap hex codes (worst #FECACA -> best #A7F3D0).
\newcommand{\HmSpeedIsoLive}{A7F3D0}
\newcommand{\HmSpeedZip}{CEF8A9}
\newcommand{\HmSpeedMoonTiny}{E3F9A7}
\newcommand{\HmSpeedVoskSmall}{E9F9AD}
\newcommand{\HmSpeedVoskLarge}{FAF9BF}
\newcommand{\HmSpeedNemoFiveSixty}{FEEBB9}
\newcommand{\HmSpeedParakeetOnnx}{EAF9AE}
\newcommand{\HmSpeedWhisperTiny}{FED3B3}
\newcommand{\HmSpeedWhisperBase}{FECACA}
\newcommand{\HmCoreIsoLive}{A7F3D0}
\newcommand{\HmCoreZip}{CEF8A8}
\newcommand{\HmCoreMoonTiny}{E3F9A7}
\newcommand{\HmCoreVoskSmall}{E9F9AD}
\newcommand{\HmCoreVoskLarge}{FAF9BF}
\newcommand{\HmCoreNemoFiveSixty}{FEEBB9}
\newcommand{\HmCoreParakeetOnnx}{EAF9AE}
\newcommand{\HmCoreWhisperTiny}{FED3B3}
\newcommand{\HmCoreWhisperBase}{FECACA}
\newcommand{\HmWerIsoLive}{C6F7B0}
\newcommand{\HmWerZip}{C2F6B5}
\newcommand{\HmWerMoonTiny}{F9F9BE}
\newcommand{\HmWerVoskSmall}{FED3B3}
\newcommand{\HmWerVoskLarge}{EDF9B2}
\newcommand{\HmWerNemoFiveSixty}{AAF3CD}
\newcommand{\HmWerParakeetOnnx}{A7F3D0}
\newcommand{\HmWerWhisperTiny}{FECACA}
\newcommand{\HmWerWhisperBase}{FED2B6}
\newcommand{\HmCleanIsoLive}{B2F4C4}
\newcommand{\HmCleanZip}{B0F4C7}
\newcommand{\HmCleanMoonTiny}{E6F9AA}
\newcommand{\HmCleanVoskSmall}{FEE0B0}
\newcommand{\HmCleanVoskLarge}{DFF9A3}
\newcommand{\HmCleanNemoFiveSixty}{A9F3CE}
\newcommand{\HmCleanParakeetOnnx}{A7F3D0}
\newcommand{\HmCleanWhisperTiny}{FECACA}
\newcommand{\HmCleanWhisperBase}{FED3B4}
\newcommand{\HmOtherIsoLive}{D9F99D}
\newcommand{\HmOtherZip}{D4F8A2}
\newcommand{\HmOtherMoonTiny}{FEECBA}
\newcommand{\HmOtherVoskSmall}{FECCC4}
\newcommand{\HmOtherVoskLarge}{FCF9C1}
\newcommand{\HmOtherNemoFiveSixty}{ABF3CC}
\newcommand{\HmOtherParakeetOnnx}{A7F3D0}
\newcommand{\HmOtherWhisperTiny}{FECACA}
\newcommand{\HmOtherWhisperBase}{FED1B8}
\newcommand{\HmRssIsoLive}{A7F3D0}
\newcommand{\HmRssZip}{B7F5C0}
\newcommand{\HmRssMoonTiny}{C4F6B3}
\newcommand{\HmRssVoskSmall}{CFF8A7}
\newcommand{\HmRssVoskLarge}{FECACA}
\newcommand{\HmRssNemoFiveSixty}{FEF7C1}
\newcommand{\HmRssParakeetOnnx}{FEE5B5}
\newcommand{\HmRssWhisperTiny}{C9F7AD}
\newcommand{\HmRssWhisperBase}{DBF99F}
\newcommand{\HmModelIsoLive}{B4F5C3}
\newcommand{\HmModelZip}{BEF6B9}
\newcommand{\HmModelMoonTiny}{A7F3D0}
\newcommand{\HmModelVoskSmall}{BDF6BA}
\newcommand{\HmModelVoskLarge}{FECACA}
\newcommand{\HmModelNemoFiveSixty}{FEE3B3}
\newcommand{\HmModelParakeetOnnx}{FEE5B4}
\newcommand{\HmModelWhisperTiny}{C1F6B6}
\newcommand{\HmModelWhisperBase}{DEF9A2}
\newcommand{\HmEnergyIsoLive}{ABF4CB}
\newcommand{\HmEnergyZip}{C7F7B0}
\newcommand{\HmEnergyMoonTiny}{E0F9A5}
\newcommand{\HmEnergyVoskSmall}{E3F9A7}
\newcommand{\HmEnergyVoskLarge}{F8F9BD}
\newcommand{\HmEnergyNemoFiveSixty}{FEEAB8}
\newcommand{\HmEnergyParakeetOnnx}{EDF9B2}
\newcommand{\HmEnergyWhisperTiny}{FED4B2}
\newcommand{\HmEnergyWhisperBase}{FECACA}


\renewcommand{\briefrunningtitle}{Executive Brief}
\hypersetup{
  pdftitle={Isosyll v4.0.0 — Executive Brief},
  pdfauthor={NEVERHUMAN Research},
  pdfsubject={An always-listening speech recognizer on one CPU core},
  pdfkeywords={Isosyll, speech recognition, streaming ASR, on-device, energy, CPU}
}

\begin{document}
\thispagestyle{briefpage1}

\briefmasthead{Executive Brief \;\textbar\; For CEOs, Product Leaders \& Operations Executives}%
  {Always Listening, Barely Computing: Speech Recognition on One CPU Core}%
  {Isosyll turns live speech into text on \textbf{\ColdRcTwoLiveLmCorePct\% of one ordinary CPU core}, with no cloud round trip, no GPU and no Python stack --- \textbf{\VsSpeedXOne\x{} the single-core throughput of the fastest open streaming rival we measured, at \VsEnergyPctOne\% of its CPU-core energy}.}

% ---- 4-Column Hero KPI Strip ------------------------------------------------
\noindent
\begin{minipage}[t]{0.242\textwidth}
  \vspace{0pt}%
  \kpicard{emerald}{\VsSpeedXOne\x}{Throughput vs.\ Fastest Open Rival}{Paired against Zipformer INT8 on the same core in the same timing blocks: \textbf{\ColdRcTwoLiveLmSpeed\x} vs.\ \ColdZipRsixSpeed\x{} real time.}
\end{minipage}\hfill
\begin{minipage}[t]{0.242\textwidth}
  \vspace{0pt}%
  \kpicard{navy}{\VsEnergyPctOne\%}{Of the Rival's Core Energy}{\textbf{\ColdRcTwoLiveLmWatts\,J} vs.\ \ColdZipRsixWatts\,J of idle-subtracted CPU-core energy per audio-second (not wall power).}
\end{minipage}\hfill
\begin{minipage}[t]{0.242\textwidth}
  \vspace{0pt}%
  \kpicard{teal}{\ColdRcTwoLiveLmCorePct\%}{Of One Core per Live Stream}{\ColdRcTwoLiveLmWer\% word error rate live in \ColdRcTwoLiveLmRss\,MB RAM (optional end-of-sentence pass reaches \ColdRcTwoFinalWer\%).}
\end{minipage}\hfill
\begin{minipage}[t]{0.242\textwidth}
  \vspace{0pt}%
  \kpicard{blue}{1 Command}{One Signed Binary}{\textasciitilde 1.6\,MB pure-Rust program for Linux (x86-64, arm64) and Apple-silicon Macs. Audio stays on-device.}
\end{minipage}

\vspace{4.5pt}

% ---- Executive Narrative ----------------------------------------------------
\noindent
\begin{minipage}[t]{0.486\textwidth}
  \vspace{0pt}
  \colheading{Always-On Speech Has a Meter That Never Stops}
  {\fontsize{8.25}{10.2}\selectfont
  A recognizer that waits for a wake word or a button press misses what matters. An \emph{always-on} recognizer --- for meeting capture, call-centre assist, in-car and wearable assistants, accessibility captions --- must process \textbf{every second of audio, whether anyone is speaking or not}. That bill never pauses: it is paid in cloud spend, in battery life, and in heat.

  \vspace{2pt}
  Streaming every microphone to a cloud GPU adds a network round trip, a per-minute bill, and a privacy question every customer and regulator asks: \textbf{where does the recording go?} Keeping speech on the device answers it --- if an ordinary CPU can afford to listen.
  }
\end{minipage}\hfill
\begin{minipage}[t]{0.486\textwidth}
  \vspace{0pt}
  \colheading{What Isosyll Changes: Listening Becomes Cheap}
  {\fontsize{8.25}{10.2}\selectfont
  \textbf{Isosyll} is a compact streaming recognizer that writes words as a person talks at \textbf{\ColdRcTwoLiveLmSpeed\x{} real time} --- an hour of speech in under two minutes --- using \ColdRcTwoLiveLmCorePct\% of one CPU core and \ColdRcTwoLiveLmRss\,MB of memory.

  \vspace{2pt}
  Against Zipformer INT8, the fastest open streaming rival we measured, Isosyll delivers \textbf{\VsSpeedXOne\x{} the single-core throughput at \VsEnergyPctOne\% of its CPU-core energy}: the same core transcribes \VsSpeedXOne\x{} as much audio for under 60\% of the energy. An optional end-of-sentence commit pass on the same \ColdRcTwoLiveLmModel\,MB weights reaches \ColdRcTwoFinalWer\% errors at \ColdRcTwoFinalSpeed\x{} when final polish is needed.
  }
\end{minipage}

\vspace{5pt}

% ---- Canonical Vector Hero Graphic (exact same figure as the paper) ---------
\begin{tcolorbox}[enhanced,colback=white,colframe=hair,boxrule=0.5pt,arc=2.5pt,
  left=6pt,right=6pt,top=5.5pt,bottom=4.5pt,before skip=0pt,after skip=0pt]
  \centering
  \briefherographic{0.94\linewidth}{4.62in}\par
  \vspace{2.5pt}
  \raggedright
  \noindent{\sffamily\fontsize{6.2}{7.7}\selectfont\color{sub}\textbf{\color{ink}Executive takeaway:} Every recognizer decoded the same 730 LibriSpeech DEV recordings (\ColdAudioMinutes\,min) on one pinned AMD Zen~2 core. Isosyll leads the streaming field on throughput, and on energy among the metered systems; on the unselected half of DEV its live text trails Zipformer by \CandLiveLmCDelta\,pp WER (bounds \CandLiveLmCLo{} to \CandLiveLmCHi).}
\end{tcolorbox}

% ==================== PAGE 2 =================================================
\clearpage
\thispagestyle{briefpage2}

\briefsection{01}{Where Isosyll Creates Immediate Advantage}

\noindent
\begin{minipage}[t]{0.324\textwidth}
  \vspace{0pt}%
  \playcard{Play 01}{emerald}{Cheaper Listening Fleets}{With \VsSpeedXOne\x{} the single-core throughput of the fastest open streaming rival, each CPU core handles \VsSpeedXOne\x{} as much call and meeting audio for under 60\% of the core energy.}
\end{minipage}\hfill
\begin{minipage}[t]{0.324\textwidth}
  \vspace{0pt}%
  \playcard{Play 02}{navy}{Private by Default, On-Device}{Speech is transcribed where it is captured --- laptops, kiosks, vehicles, edge boxes. No cloud round trip, no per-minute API bill, and recordings never leave the building.}
\end{minipage}\hfill
\begin{minipage}[t]{0.324\textwidth}
  \vspace{0pt}%
  \playcard{Play 03}{teal}{Zero-GPU Power Profile}{At \ColdRcTwoLiveLmWatts\,J of CPU-core energy per audio-second (about \ColdRcTwoLiveLmWatts\,W for one real-time stream) and \ColdRcTwoLiveLmRss\,MB RAM, always-on captions run on ordinary CPUs, no GPU accelerator needed.}
\end{minipage}

\vspace{2.5pt}
\briefsection{02}{Measured Scoreboard at a Glance (Directional Heat Map)}

\noindent{\fontsize{7.8}{9.6}\selectfont\color{sub}Same 730 LibriSpeech DEV recordings, one pinned Zen~2 CPU core, one scorer, median of three cold runs. \textbf{\color{ink}Cell heat map:} \hmlegend{} emerald = best in column, coral = worst (log-scaled for speed, memory and energy). Showing our single fastest streaming model vs.\ public rivals.}

\vspace{2.5pt}
\noindent
{\fontsize{7.4}{9.2}\selectfont\setlength{\tabcolsep}{4pt}\renewcommand{\arraystretch}{1.20}%
\begin{tabularx}{\textwidth}{@{}L{4.55cm}R{1.55cm}R{1.95cm}R{1.75cm}R{1.50cm}Y@{}}
\toprule
\thd{Recognizer (mode)} & \thd{Speed $\uparrow$} & \thd{Word errors $\downarrow$} & \thd{Core J / s $\downarrow$} & \thd{Peak RAM $\downarrow$} & \thd{What it means} \\
\midrule
\textbf{Isosyll} (streaming, 160\,ms) & \hmcell{\HmSpeedIsoLive}{\textbf{\ColdRcTwoLiveLmSpeed\x}} & \hmcell{\HmWerIsoLive}{\textbf{\ColdRcTwoLiveLmWer\%}} & \hmcell{\HmEnergyIsoLive}{\textbf{\ColdRcTwoLiveLmWatts\,J$^\ast$}} & \hmcell{\HmRssIsoLive}{\textbf{\ColdRcTwoLiveLmRss\,MB}} & \textbf{\VsSpeedXOne\x{} Zipformer's speed, \VsEnergyPctOne\% of its energy} \\
Zipformer INT8 (streaming) & \hmcell{\HmSpeedZip}{\ColdZipRsixSpeed\x} & \hmcell{\HmWerZip}{\ColdZipRsixWer\%} & \hmcell{\HmEnergyZip}{\ColdZipRsixWatts\,J$^\ast$} & \hmcell{\HmRssZip}{\ColdZipRsixRss\,MB} & Fastest open streaming rival \\
Moonshine Tiny (streaming) & \hmcell{\HmSpeedMoonTiny}{\rone{\ColdMoonTinySpeed}\x} & \hmcell{\HmWerMoonTiny}{\ColdMoonTinyWer\%} & \hmna & \hmcell{\HmRssMoonTiny}{\ColdMoonTinyRss\,MB} & Slower, with more errors \\
Vosk Small (streaming) & \hmcell{\HmSpeedVoskSmall}{\rone{\ColdVoskSmallSpeed}\x} & \hmcell{\HmWerVoskSmall}{\ColdVoskSmallWer\%} & \hmna & \hmcell{\HmRssVoskSmall}{\ColdVoskSmallRss\,MB} & Slower and less accurate \\
Vosk Large\S{} (streaming) & \hmcell{\HmSpeedVoskLarge}{\rone{\ColdVoskLargeSpeed}\x} & \hmcell{\HmWerVoskLarge}{\ColdVoskLargeWer\%} & \hmna & \hmcell{\HmRssVoskLarge}{\ColdVoskLargeRss\,MB} & Heavy memory footprint (\textgreater 7\,GB) \\
Nemotron 0.6B, 560\,ms\S{} (streaming) & \hmcell{\HmSpeedNemoFiveSixty}{\rone{\ColdNemoFiveSixtySpeed}\x} & \hmcell{\HmWerNemoFiveSixty}{\ColdNemoFiveSixtyWer\%} & \hmna & \hmcell{\HmRssNemoFiveSixty}{\ColdNemoFiveSixtyRss\,MB} & 0.6B params: more accurate, far slower \\
Parakeet TDT v3 (offline) & \hmcell{\HmSpeedParakeetOnnx}{\rone{\ColdParakeetOnnxAsrSpeed}\x} & \hmcell{\HmWerParakeetOnnx}{\ColdParakeetOnnxAsrWer\%} & \hmna & \hmcell{\HmRssParakeetOnnx}{\ColdParakeetOnnxAsrRss\,MB} & More accurate; needs whole clip first \\
Whisper.cpp Tiny\S{} (streaming) & \hmcell{\HmSpeedWhisperTiny}{\rone{\ColdWhisperCppTinySpeed}\x} & \hmcell{\HmWerWhisperTiny}{\ColdWhisperCppTinyWer\%} & \hmna & \hmcell{\HmRssWhisperTiny}{\ColdWhisperCppTinyRss\,MB} & Cannot keep up with live audio \\
Whisper.cpp Base\S{} (streaming) & \hmcell{\HmSpeedWhisperBase}{\rone{\ColdWhisperCppBaseSpeed}\x} & \hmcell{\HmWerWhisperBase}{\ColdWhisperCppBaseWer\%} & \hmna & \hmcell{\HmRssWhisperBase}{\ColdWhisperCppBaseRss\,MB} & Cannot keep up with live audio \\
\bottomrule
\end{tabularx}}
\vspace{3pt}
\noindent{\sffamily\fontsize{6.15}{7.6}\selectfont\color{sub}$^\ast$\,Idle-subtracted CPU-core joules per audio-second (not wall power), timed side by side in the same interleaved blocks; other rivals ran in separate sessions so their energy is not shown here. \S\,Speed timed on a 73-clip subset; WER on all 730. Optional utterance-end commit pass on the same \ColdRcTwoLiveLmModel\,MB weights: \ColdRcTwoFinalWer\% WER at \ColdRcTwoFinalSpeed\x{} (\ColdRcTwoFinalWatts\,J, \ColdRcTwoFinalRss\,MB). Receipts: \texttt{docs/paper/data/}.}

\vspace{3.5pt}
\briefsection{03}{How to Pilot in 30 Days --- and What Leadership Should Know}

\noindent
\begin{minipage}[t]{0.478\textwidth}
  \vspace{0pt}
  \colheading{3-Step Low-Risk Pilot}
  \vspace{1pt}
  {\fontsize{7.75}{9.7}\selectfont\raggedright
  \noindent\stepbadge{1}\;\textbf{Install and try it today (Day 1).}\\
  One command installs a signed, checksummed binary; a second transcribes a recording.
  }
  \begin{codecard}{Install \& Run (Linux x86-64 / arm64, Apple-Silicon Mac)}
curl -fsSL https://github.com/neverhuman/jsr/releases/\textbackslash\\
\hspace*{1.2em}latest/download/install.sh | sh\\
isosyll transcribe meeting.wav
  \end{codecard}
  {\fontsize{7.75}{9.7}\selectfont\raggedright

  \vspace{2.5pt}
  \noindent\stepbadge{2}\;\textbf{Measure on your own audio and hardware (Weeks 1--2).}\\
  Run it on a sample of your calls, meetings or field recordings on machines you already own. Compare word errors and CPU load with your current vendor.

  \vspace{2.5pt}
  \noindent\stepbadge{3}\;\textbf{Size the business case (Week 3).}\\
  Convert measured streams per core into fleet size, cloud spend avoided and battery hours. Decide where on-device listening replaces a cloud API call.
  }

  \vspace{3pt}
  \begin{execbox}[emerald]{Why Leaders Can Trust the Numbers}
    Every figure comes from published measurement records. Isosyll produces \textbf{byte-identical transcripts} to the signed reference decoders on all \ScalarUtts{} panel recordings, on both Intel/AMD and Apple-silicon processors. Releases are SSH-signed and checksummed.
  \end{execbox}
\end{minipage}\hfill
\begin{minipage}[t]{0.495\textwidth}
  \vspace{0pt}
  \begin{govbox}{What This Is Not --- Honest Limits}
    \textbf{Release \releasever{} is strong evidence, not a guarantee. Six limits to keep in view:}
    \begin{itemize}[leftmargin=1.10em,itemsep=1.5pt,topsep=2pt]
      \item \textbf{Exploratory, released by decision:} Isosyll \releasever{} \textbf{missed its pre-registered full-DEV edit gates} and was released by owner decision; all results are exploratory.
      \item \textbf{One benchmark, read English:} All accuracy is LibriSpeech development data (audiobooks). No test-set results are claimed; conversational, accented or noisy speech may rank differently.
      \item \textbf{Live text trails slightly:} On the unselected half of DEV, live Isosyll makes \CandLiveLmCDelta\,pp more word errors than Zipformer (bounds \CandLiveLmCLo{} to \CandLiveLmCHi). The optional end-of-sentence pass closes the gap (\ColdRcTwoFinalWer\% WER).
      \item \textbf{Energy edge is live mode:} the optional commit pass uses \VsEnergyPctTwo\% of Zipformer's CPU-core energy and \ColdRcTwoFinalRss\,MB RAM.
      \item \textbf{Bigger models are more accurate:} 0.6B-parameter models (Nemotron, Parakeet) make fewer errors at far higher compute.
      \item \textbf{Power scope:} Measured energy is CPU-core joules only (not wall/battery power); GPU TDP figures (70--250\,W) are external hardware context, not a benchmark measurement. No latency claim is made here.
    \end{itemize}
  \end{govbox}
\end{minipage}

\brieffooterstrip{Ready to Try Isosyll or Share It With Your Engineering Team?}%
  {\textbf{Next reads:} CTO Technical \& Architecture Brief (\texttt{docs/briefs/isosyll-v4-cto-brief.pdf}) \;\textbar\; Full Paper (\texttt{docs/paper/paper.pdf}) \;\textbar\; Install Guide (\texttt{docs/install.md}) \;\textbar\; Release Records (\texttt{docs/ultra/releases/})}%
  {https://github.com/neverhuman/jsr}

\end{document}
